\IfFileExists{stacks-project.cls}{%
\documentclass{stacks-project}
}{%
\documentclass{amsart}
}

% For dealing with references we use the comment environment
\usepackage{verbatim}
\newenvironment{reference}{\comment}{\endcomment}
%\newenvironment{reference}{}{}
\newenvironment{slogan}{\comment}{\endcomment}
\newenvironment{history}{\comment}{\endcomment}

% For commutative diagrams we use Xy-pic
\usepackage[all]{xy}

% We use 2cell for 2-commutative diagrams.
\xyoption{2cell}
\UseAllTwocells

% We use multicol for the list of chapters between chapters
\usepackage{multicol}

% This is generally recommended for better output
\usepackage{lmodern}
\usepackage[T1]{fontenc}

% For cross-file-references
\usepackage{xr-hyper}

% Package for hypertext links:
\usepackage{hyperref}

% For any local file, say "hello.tex" you want to link to please
% use \externaldocument[hello-]{hello}
\externaldocument[introduction-]{introduction}
\externaldocument[conventions-]{conventions}
\externaldocument[sets-]{sets}
\externaldocument[categories-]{categories}
\externaldocument[topology-]{topology}
\externaldocument[sheaves-]{sheaves}
\externaldocument[sites-]{sites}
\externaldocument[stacks-]{stacks}
\externaldocument[fields-]{fields}
\externaldocument[algebra-]{algebra}
\externaldocument[brauer-]{brauer}
\externaldocument[homology-]{homology}
\externaldocument[derived-]{derived}
\externaldocument[simplicial-]{simplicial}
\externaldocument[more-algebra-]{more-algebra}
\externaldocument[smoothing-]{smoothing}
\externaldocument[modules-]{modules}
\externaldocument[sites-modules-]{sites-modules}
\externaldocument[injectives-]{injectives}
\externaldocument[cohomology-]{cohomology}
\externaldocument[sites-cohomology-]{sites-cohomology}
\externaldocument[dga-]{dga}
\externaldocument[dpa-]{dpa}
\externaldocument[sdga-]{sdga}
\externaldocument[hypercovering-]{hypercovering}
\externaldocument[schemes-]{schemes}
\externaldocument[constructions-]{constructions}
\externaldocument[properties-]{properties}
\externaldocument[morphisms-]{morphisms}
\externaldocument[coherent-]{coherent}
\externaldocument[divisors-]{divisors}
\externaldocument[limits-]{limits}
\externaldocument[varieties-]{varieties}
\externaldocument[topologies-]{topologies}
\externaldocument[descent-]{descent}
\externaldocument[perfect-]{perfect}
\externaldocument[more-morphisms-]{more-morphisms}
\externaldocument[flat-]{flat}
\externaldocument[groupoids-]{groupoids}
\externaldocument[more-groupoids-]{more-groupoids}
\externaldocument[etale-]{etale}
\externaldocument[chow-]{chow}
\externaldocument[intersection-]{intersection}
\externaldocument[pic-]{pic}
\externaldocument[weil-]{weil}
\externaldocument[adequate-]{adequate}
\externaldocument[dualizing-]{dualizing}
\externaldocument[duality-]{duality}
\externaldocument[discriminant-]{discriminant}
\externaldocument[derham-]{derham}
\externaldocument[local-cohomology-]{local-cohomology}
\externaldocument[algebraization-]{algebraization}
\externaldocument[curves-]{curves}
\externaldocument[resolve-]{resolve}
\externaldocument[models-]{models}
\externaldocument[functors-]{functors}
\externaldocument[equiv-]{equiv}
\externaldocument[pione-]{pione}
\externaldocument[etale-cohomology-]{etale-cohomology}
\externaldocument[proetale-]{proetale}
\externaldocument[relative-cycles-]{relative-cycles}
\externaldocument[more-etale-]{more-etale}
\externaldocument[trace-]{trace}
\externaldocument[crystalline-]{crystalline}
\externaldocument[spaces-]{spaces}
\externaldocument[spaces-properties-]{spaces-properties}
\externaldocument[spaces-morphisms-]{spaces-morphisms}
\externaldocument[decent-spaces-]{decent-spaces}
\externaldocument[spaces-cohomology-]{spaces-cohomology}
\externaldocument[spaces-limits-]{spaces-limits}
\externaldocument[spaces-divisors-]{spaces-divisors}
\externaldocument[spaces-over-fields-]{spaces-over-fields}
\externaldocument[spaces-topologies-]{spaces-topologies}
\externaldocument[spaces-descent-]{spaces-descent}
\externaldocument[spaces-perfect-]{spaces-perfect}
\externaldocument[spaces-more-morphisms-]{spaces-more-morphisms}
\externaldocument[spaces-flat-]{spaces-flat}
\externaldocument[spaces-groupoids-]{spaces-groupoids}
\externaldocument[spaces-more-groupoids-]{spaces-more-groupoids}
\externaldocument[bootstrap-]{bootstrap}
\externaldocument[spaces-pushouts-]{spaces-pushouts}
\externaldocument[spaces-chow-]{spaces-chow}
\externaldocument[groupoids-quotients-]{groupoids-quotients}
\externaldocument[spaces-more-cohomology-]{spaces-more-cohomology}
\externaldocument[spaces-simplicial-]{spaces-simplicial}
\externaldocument[spaces-duality-]{spaces-duality}
\externaldocument[formal-spaces-]{formal-spaces}
\externaldocument[restricted-]{restricted}
\externaldocument[spaces-resolve-]{spaces-resolve}
\externaldocument[formal-defos-]{formal-defos}
\externaldocument[defos-]{defos}
\externaldocument[cotangent-]{cotangent}
\externaldocument[examples-defos-]{examples-defos}
\externaldocument[algebraic-]{algebraic}
\externaldocument[examples-stacks-]{examples-stacks}
\externaldocument[stacks-sheaves-]{stacks-sheaves}
\externaldocument[criteria-]{criteria}
\externaldocument[artin-]{artin}
\externaldocument[quot-]{quot}
\externaldocument[stacks-properties-]{stacks-properties}
\externaldocument[stacks-morphisms-]{stacks-morphisms}
\externaldocument[stacks-limits-]{stacks-limits}
\externaldocument[stacks-cohomology-]{stacks-cohomology}
\externaldocument[stacks-perfect-]{stacks-perfect}
\externaldocument[stacks-introduction-]{stacks-introduction}
\externaldocument[stacks-more-morphisms-]{stacks-more-morphisms}
\externaldocument[stacks-geometry-]{stacks-geometry}
\externaldocument[moduli-]{moduli}
\externaldocument[moduli-curves-]{moduli-curves}
\externaldocument[examples-]{examples}
\externaldocument[exercises-]{exercises}
\externaldocument[guide-]{guide}
\externaldocument[desirables-]{desirables}
\externaldocument[coding-]{coding}
\externaldocument[obsolete-]{obsolete}
\externaldocument[fdl-]{fdl}
\externaldocument[index-]{index}

% Theorem environments.
%
\theoremstyle{plain}
\newtheorem{theorem}[subsection]{Theorem}
\newtheorem{proposition}[subsection]{Proposition}
\newtheorem{lemma}[subsection]{Lemma}

\theoremstyle{definition}
\newtheorem{definition}[subsection]{Definition}
\newtheorem{example}[subsection]{Example}
\newtheorem{exercise}[subsection]{Exercise}
\newtheorem{situation}[subsection]{Situation}

\theoremstyle{remark}
\newtheorem{remark}[subsection]{Remark}
\newtheorem{remarks}[subsection]{Remarks}

\numberwithin{equation}{subsection}

% Macros
%
\def\lim{\mathop{\mathrm{lim}}\nolimits}
\def\colim{\mathop{\mathrm{colim}}\nolimits}
\def\Spec{\mathop{\mathrm{Spec}}}
\def\Hom{\mathop{\mathrm{Hom}}\nolimits}
\def\Ext{\mathop{\mathrm{Ext}}\nolimits}
\def\SheafHom{\mathop{\mathcal{H}\!\mathit{om}}\nolimits}
\def\SheafExt{\mathop{\mathcal{E}\!\mathit{xt}}\nolimits}
\def\Sch{\mathit{Sch}}
\def\Mor{\mathop{\mathrm{Mor}}\nolimits}
\def\Ob{\mathop{\mathrm{Ob}}\nolimits}
\def\Sh{\mathop{\mathit{Sh}}\nolimits}
\def\NL{\mathop{N\!L}\nolimits}
\def\CH{\mathop{\mathrm{CH}}\nolimits}
\def\proetale{{pro\text{-}\acute{e}tale}}
\def\etale{{\acute{e}tale}}
\def\QCoh{\mathit{QCoh}}
\def\Ker{\mathop{\mathrm{Ker}}}
\def\Im{\mathop{\mathrm{Im}}}
\def\Coker{\mathop{\mathrm{Coker}}}
\def\Coim{\mathop{\mathrm{Coim}}}

% Boxtimes
%
\DeclareMathSymbol{\boxtimes}{\mathbin}{AMSa}{"02}

%
% Macros for moduli stacks/spaces
%
\def\QCohstack{\mathcal{QC}\!\mathit{oh}}
\def\Cohstack{\mathcal{C}\!\mathit{oh}}
\def\Spacesstack{\mathcal{S}\!\mathit{paces}}
\def\Quotfunctor{\mathrm{Quot}}
\def\Hilbfunctor{\mathrm{Hilb}}
\def\Curvesstack{\mathcal{C}\!\mathit{urves}}
\def\Polarizedstack{\mathcal{P}\!\mathit{olarized}}
\def\Complexesstack{\mathcal{C}\!\mathit{omplexes}}
% \Pic is the operator that assigns to X its picard group, usage \Pic(X)
% \Picardstack_{X/B} denotes the Picard stack of X over B
% \Picardfunctor_{X/B} denotes the Picard functor of X over B
\def\Pic{\mathop{\mathrm{Pic}}\nolimits}
\def\Picardstack{\mathcal{P}\!\mathit{ic}}
\def\Picardfunctor{\mathrm{Pic}}
\def\Deformationcategory{\mathcal{D}\!\mathit{ef}}
\usepackage{cleveref}
\crefname{originallemma}{original lemma}{original lemmas}
\Crefname{originallemma}{Original lemma}{Original lemmas}
\makeatletter
\newlabel{lemma-improve-distinguished-triangle-homotopy}{{0G6C}{source}{}{sourcepointer-1}{}}
\newlabel{lemma-improve-distinguished-triangle-homotopy@cref}{{[originallemma][0][]0G6C}{[1][0][]source}{}{}{}}
\makeatother
\crefname{originallemma}{original lemma}{original lemmas}
\Crefname{originallemma}{Original lemma}{Original lemmas}
\makeatletter
\newlabel{lemma-subcategory-left-resolution}{{05T7}{source}{}{sourcepointer-2}{}}
\newlabel{lemma-subcategory-left-resolution@cref}{{[originallemma][0][]05T7}{[1][0][]source}{}{}{}}
\makeatother
\crefname{originalproposition}{original proposition}{original propositions}
\Crefname{originalproposition}{Original proposition}{Original propositions}
\makeatletter
\newlabel{proposition-homotopy-category-triangulated}{{014S}{source}{}{sourcepointer-3}{}}
\newlabel{proposition-homotopy-category-triangulated@cref}{{[originalproposition][0][]014S}{[1][0][]source}{}{}{}}
\makeatother
\title{Special direct systems built from left resolutions}
\author{The Stacks Project Authors}
\date{stacks-project:derived.tex:local-snapshot-20261012; GNU FDL 1.2 or later}
\begin{document}
\maketitle
\noindent\textit{Allowed standard prerequisites (editorial).}\begin{itemize}
\item Left resolution of complexes. lemma-subcategory-left-resolution (same chapter, carded in the old generation). Supplies the left decomposition engine applied twice in the inductive construction.
\item Improving triangle morphisms up to homotopy. lemma-improve-distinguished-triangle-homotopy (same chapter, carded in the old generation). Supplies the homotopy-improvement engine used to glue the triangle at each stage.
\item The homotopy category is triangulated. proposition-homotopy-category-triangulated (same chapter, carded in the old generation). Supplies the triangulated-category bookkeeping used to assemble the special direct system.
\end{itemize}\medskip
\ifdefined\hypertarget\hypertarget{sourcepointer-1}{}\fi
\noindent\textit{Source reference.} Original lemma 0G6C: sources/stacks-derived/pinned-original-source.tex, source lines 3676--3676. Navigation pointer only: the counted proof cites this statement by \textbackslash{}ref as a standing input; the deduction the counted proof has to make is carried out inside the counted span, and no part of that deduction is delegated to this pointer.\par\medskip
\ifdefined\hypertarget\hypertarget{sourcepointer-2}{}\fi
\noindent\textit{Source reference.} Original lemma 05T7: sources/stacks-derived/pinned-original-source.tex, source lines 5519--5519. Navigation pointer only: the counted proof cites this statement by \textbackslash{}ref as a standing input; the deduction the counted proof has to make is carried out inside the counted span, and no part of that deduction is delegated to this pointer.\par\medskip
\ifdefined\hypertarget\hypertarget{sourcepointer-3}{}\fi
\noindent\textit{Source reference.} Original proposition 014S: sources/stacks-derived/pinned-original-source.tex, source lines 3537--3537. Navigation pointer only: the counted proof cites this statement by \textbackslash{}ref as a standing input; the deduction the counted proof has to make is carried out inside the counted span, and no part of that deduction is delegated to this pointer.\par\medskip
\begin{lemma}
\label{lemma-special-direct-system}
Let $\mathcal{A}$ be an abelian category. Let
$\mathcal{P} \subset \Ob(\mathcal{A})$ be a subset.
Assume $\mathcal{P}$ contains $0$, is closed under (finite) direct sums,
and every object of $\mathcal{A}$ is a quotient of an
element of $\mathcal{P}$. Let $K^\bullet$ be a complex.
There exists a commutative diagram
$$
\xymatrix{
P_1^\bullet \ar[d] \ar[r] & P_2^\bullet \ar[d] \ar[r] & \ldots \\
\tau_{\leq 1}K^\bullet \ar[r] & \tau_{\leq 2}K^\bullet \ar[r] & \ldots
}
$$
in the category of complexes such that
\begin{enumerate}
\item the vertical arrows are quasi-isomorphisms and termwise surjective,
\item $P_n^\bullet$ is a bounded above complex with terms in
$\mathcal{P}$,
\item the arrows $P_n^\bullet \to P_{n + 1}^\bullet$
are termwise split injections and each cokernel
$P^i_{n + 1}/P^i_n$ is an element of $\mathcal{P}$.
\end{enumerate}
\end{lemma}

\begin{proof}
We are going to use that the homotopy category $K(\mathcal{A})$ is a
triangulated category, see Proposition
\ref{proposition-homotopy-category-triangulated}.
By Lemma \ref{lemma-subcategory-left-resolution} we can find a
termwise surjective map of complexes $P_1^\bullet \to \tau_{\leq 1}K^\bullet$
which is a quasi-isomorphism
such that the terms of $P_1^\bullet$ are in $\mathcal{P}$.
By induction it suffices, given
$P_1^\bullet, \ldots, P_n^\bullet$ to construct
$P_{n + 1}^\bullet$ and the maps
$P_n^\bullet \to P_{n + 1}^\bullet$ and
$P_{n + 1}^\bullet \to \tau_{\leq n + 1}K^\bullet$.

\medskip\noindent
Choose a distinguished triangle
$P_n^\bullet \to \tau_{\leq n + 1}K^\bullet \to C^\bullet \to P_n^\bullet[1]$
in $K(\mathcal{A})$. Applying
Lemma \ref{lemma-subcategory-left-resolution} we choose a
map of complexes $Q^\bullet \to C^\bullet$
which is a quasi-isomorphism such that the terms of $Q^\bullet$
are in $\mathcal{P}$. By the axioms of triangulated categories
we may fit the composition $Q^\bullet \to C^\bullet \to P_n^\bullet[1]$ into
a distinguished triangle
$P_n^\bullet \to P_{n + 1}^\bullet \to Q^\bullet \to P_n^\bullet[1]$
in $K(\mathcal{A})$.
By Lemma \ref{lemma-improve-distinguished-triangle-homotopy}
we may and do assume
$0 \to P_n^\bullet \to P_{n + 1}^\bullet \to Q^\bullet \to 0$
is a termwise split short exact sequence. This implies that
the terms of $P_{n + 1}^\bullet$ are in $\mathcal{P}$ and that
$P_n^\bullet \to P_{n + 1}^\bullet$ is a termwise split injection
whose cokernels are in $\mathcal{P}$.
By the axioms of triangulated categories we obtain a map
of distinguished triangles
$$
\xymatrix{
P_n^\bullet \ar[r] \ar[d] &
P_{n + 1}^\bullet \ar[r] \ar[d] &
Q^\bullet \ar[r] \ar[d] &
P_n^\bullet[1] \ar[d] \\
P_n^\bullet \ar[r] &
\tau_{\leq n + 1}K^\bullet \ar[r] &
C^\bullet \ar[r] &
P_n^\bullet[1]
}
$$
in the triangulated category $K(\mathcal{A})$. Choose an actual morphism of
complexes $f : P_{n + 1}^\bullet \to \tau_{\leq n + 1}K^\bullet$.
The left square of the diagram above commutes up to homotopy, but as
$P_n^\bullet \to P_{n + 1}^\bullet$ is a termwise split injection
we can lift the homotopy and modify our choice of $f$ to make it commute.
Finally, $f$ is a quasi-isomorphism, because both $P_n^\bullet \to P_n^\bullet$
and $Q^\bullet \to C^\bullet$ are.

\medskip\noindent
At this point we have all the properties we want, except we don't know
that the map $f : P_{n + 1}^\bullet \to \tau_{\leq n + 1}K^\bullet$
is termwise surjective. Since we have the commutative diagram
$$
\xymatrix{
P_n^\bullet \ar[d] \ar[r] & P_{n + 1}^\bullet \ar[d] \\
\tau_{\leq n}K^\bullet \ar[r] & \tau_{\leq n + 1}K^\bullet
}
$$
of complexes, by induction hypothesis we see that $f$ is surjective
on terms in all degrees except possibly $n$ and $n + 1$. Choose
an object $P \in \mathcal{P}$ and a surjection $q : P \to K^n$.
Consider the map
$$
g :
P^\bullet = (\ldots \to 0 \to P \xrightarrow{1} P \to 0 \to \ldots)
\longrightarrow
\tau_{\leq n + 1}K^\bullet
$$
with first copy of $P$ in degree $n$ and maps given by
$q$ in degree $n$ and $d_K \circ q$ in degree $n + 1$.
This is a surjection in degree $n$ and the cokernel in
degree $n + 1$ is $H^{n + 1}(\tau_{\leq n + 1}K^\bullet)$;
to see this recall that $\tau_{\leq n + 1}K^\bullet$ has
$\Ker(d_K^{n + 1})$ in degree $n + 1$.
However, since $f$ is a quasi-isomorphism we know that
$H^{n + 1}(f)$ is surjective. Hence after replacing 
$f : P_{n + 1}^\bullet \to \tau_{\leq n + 1}K^\bullet$
by
$f \oplus g : P_{n + 1}^\bullet \oplus P^\bullet \to \tau_{\leq n + 1}K^\bullet$
we win.
\end{proof}

\end{document}
